% !TEX program = xelatex
\documentclass[a4paper,11pt,openany,fontset=none]{ctexrep}
\newcommand{\reporttitle}{2026年末翻倍空间筛选}
\newcommand{\reportsubtitle}{A股高弹性标的：基准价值、牛市尾部与证据门槛}
\newcommand{\reportkicker}{A股策略与多标的研究}
\newcommand{\reportscope}{中国A股 | 当前价格 | 2026年末事件窗口}
\newcommand{\reportdate}{2026年7月22日}
\newcommand{\reportdatacutoff}{行情截至2026年7月22日收盘；财务截至2025年报、2026Q1及已披露2026H1预告}
\newcommand{\reporttype}{全市场高弹性筛选}
\newcommand{\reportauthor}{AStock研究体系}
\newcommand{\reporthouseview}{没有基准情景下高置信度翻倍标的。雅化16.79元：B+级条件核心，正式区间约25--26元，42元外部锚期限未披露；H2利润与OCF验证，现金仍负或库存/套保主导则失效。恩捷47.84元：C级零权敏感性，无正式目标；103元外部锚期限未披露，公司级量价、利用率与良率缺失即维持阻断。}
\newcommand{\reportquality}{核心价格为实时收盘数据；缺少显式点目标者零外部权重；显式目标但期限未披露者仅作限权牛市参考，且须先通过公司级估值资格门槛。}
\newcommand{\reportdisclaimer}{本报告基于公开资料整理，不构成任何证券买卖建议。}
% Reusable IB-style preamble for XeLaTeX equity research reports
% Usage: % Reusable IB-style preamble for XeLaTeX equity research reports
% Usage: % Reusable IB-style preamble for XeLaTeX equity research reports
% Usage: \input{preamble} after \documentclass[a4paper,11pt,openany,fontset=none]{ctexrep}

% ============ 页面布局 ============
\usepackage[top=2cm, bottom=2cm, left=2cm, right=2cm]{geometry}
\setlength{\parskip}{0.3em}
\setlength{\parindent}{2em}
\renewcommand{\baselinestretch}{1.15}

% ============ 字体（macOS） ============
\setCJKmainfont[BoldFont={Heiti SC}]{STSong}
\setCJKsansfont{Heiti SC}
\setCJKmonofont{Heiti SC}
\setmainfont{Times New Roman}

% ============ 颜色（IB Navy/Grey palette） ============
\usepackage{xcolor}
\definecolor{navy}{HTML}{003366}
\definecolor{steelgrey}{HTML}{4A5568}
\definecolor{lightgrey}{HTML}{F7FAFC}
\definecolor{riskred}{HTML}{C53030}
\definecolor{riskamber}{HTML}{D69E2E}
\definecolor{riskgreen}{HTML}{38A169}
\definecolor{nvgreen}{HTML}{76B900}

% ============ 表格 ============
\usepackage{booktabs}
\usepackage{longtable}
\usepackage{tabularx}
\usepackage{multirow}
\usepackage{array}
\newcolumntype{Y}{>{\centering\arraybackslash}X}
\newcolumntype{L}[1]{>{\raggedright\arraybackslash}p{#1}}
\newcolumntype{C}[1]{>{\centering\arraybackslash}p{#1}}
\newcolumntype{R}[1]{>{\raggedleft\arraybackslash}p{#1}}

% ============ 图形 ============
\usepackage{tikz}
\usetikzlibrary{shapes,arrows.meta,positioning,calc,backgrounds,fit}
\usepackage{pgfplots}
\usepgfplotslibrary{polar}
\pgfplotsset{compat=1.18}
\usepackage[edges]{forest}

% ============ 页眉页脚（报告标题和日期需在main.tex中覆盖） ============
\usepackage{fancyhdr}
\pagestyle{fancy}
\fancyhf{}
\fancyhead[L]{\small\color{steelgrey} \reporttitle}
\fancyhead[R]{\small\color{steelgrey} \reportdate}
\fancyfoot[C]{\small\color{steelgrey} \thepage}
\fancyfoot[R]{\tiny\color{steelgrey} 本报告不构成任何投资建议}
\renewcommand{\headrulewidth}{0.4pt}
\renewcommand{\footrulewidth}{0.2pt}

% ============ 章节格式 ============
\usepackage{titlesec}
\titleformat{\chapter}[display]
  {\normalfont\huge\bfseries\color{navy}}
  {\chaptertitlename\ \thechapter}{20pt}{\Huge}
\titleformat{\section}
  {\normalfont\Large\bfseries\color{navy}}
  {\thesection}{1em}{}
\titleformat{\subsection}
  {\normalfont\large\bfseries\color{steelgrey}}
  {\thesubsection}{1em}{}

% ============ 超链接 ============
\usepackage{hyperref}
\hypersetup{
  colorlinks=true,
  linkcolor=navy,
  urlcolor=navy,
  citecolor=navy,
}

% ============ 其他工具 ============
\usepackage{enumitem}
\setlist{nosep, leftmargin=2em}
\usepackage{fontawesome5}
\usepackage{amssymb}
\usepackage{pifont}
\usepackage{tcolorbox}
\tcbuselibrary{skins,breakable}
\usepackage{caption}
\captionsetup{font=small, skip=4pt}
\usepackage{float}
\setlength{\textfloatsep}{8pt plus 2pt minus 2pt}
\setlength{\floatsep}{6pt plus 2pt minus 2pt}
\setlength{\intextsep}{6pt plus 2pt minus 2pt}
\setlength{\abovecaptionskip}{4pt}
\setlength{\belowcaptionskip}{2pt}

% ============ 自定义命令 ============
\newcommand{\rmark}{\textcolor{riskred}{\ding{108}}}
\newcommand{\amark}{\textcolor{riskamber}{\ding{108}}}
\newcommand{\gmark}{\textcolor{riskgreen}{\ding{108}}}
\newcommand{\starfive}{$\bigstar\bigstar\bigstar\bigstar\bigstar$}
\newcommand{\starfour}{$\bigstar\bigstar\bigstar\bigstar$}
\newcommand{\starthree}{$\bigstar\bigstar\bigstar$}

% 信息框
\newtcolorbox{keyinsight}[1][]{
  colback=lightgrey, colframe=navy,
  fonttitle=\bfseries, title={#1}, breakable
}
\newtcolorbox{riskbox}[1][]{
  colback=red!5, colframe=riskred,
  fonttitle=\bfseries, title={#1}, breakable
}
 after \documentclass[a4paper,11pt,openany,fontset=none]{ctexrep}

% ============ 页面布局 ============
\usepackage[top=2cm, bottom=2cm, left=2cm, right=2cm]{geometry}
\setlength{\parskip}{0.3em}
\setlength{\parindent}{2em}
\renewcommand{\baselinestretch}{1.15}

% ============ 字体（macOS） ============
\setCJKmainfont[BoldFont={Heiti SC}]{STSong}
\setCJKsansfont{Heiti SC}
\setCJKmonofont{Heiti SC}
\setmainfont{Times New Roman}

% ============ 颜色（IB Navy/Grey palette） ============
\usepackage{xcolor}
\definecolor{navy}{HTML}{003366}
\definecolor{steelgrey}{HTML}{4A5568}
\definecolor{lightgrey}{HTML}{F7FAFC}
\definecolor{riskred}{HTML}{C53030}
\definecolor{riskamber}{HTML}{D69E2E}
\definecolor{riskgreen}{HTML}{38A169}
\definecolor{nvgreen}{HTML}{76B900}

% ============ 表格 ============
\usepackage{booktabs}
\usepackage{longtable}
\usepackage{tabularx}
\usepackage{multirow}
\usepackage{array}
\newcolumntype{Y}{>{\centering\arraybackslash}X}
\newcolumntype{L}[1]{>{\raggedright\arraybackslash}p{#1}}
\newcolumntype{C}[1]{>{\centering\arraybackslash}p{#1}}
\newcolumntype{R}[1]{>{\raggedleft\arraybackslash}p{#1}}

% ============ 图形 ============
\usepackage{tikz}
\usetikzlibrary{shapes,arrows.meta,positioning,calc,backgrounds,fit}
\usepackage{pgfplots}
\usepgfplotslibrary{polar}
\pgfplotsset{compat=1.18}
\usepackage[edges]{forest}

% ============ 页眉页脚（报告标题和日期需在main.tex中覆盖） ============
\usepackage{fancyhdr}
\pagestyle{fancy}
\fancyhf{}
\fancyhead[L]{\small\color{steelgrey} \reporttitle}
\fancyhead[R]{\small\color{steelgrey} \reportdate}
\fancyfoot[C]{\small\color{steelgrey} \thepage}
\fancyfoot[R]{\tiny\color{steelgrey} 本报告不构成任何投资建议}
\renewcommand{\headrulewidth}{0.4pt}
\renewcommand{\footrulewidth}{0.2pt}

% ============ 章节格式 ============
\usepackage{titlesec}
\titleformat{\chapter}[display]
  {\normalfont\huge\bfseries\color{navy}}
  {\chaptertitlename\ \thechapter}{20pt}{\Huge}
\titleformat{\section}
  {\normalfont\Large\bfseries\color{navy}}
  {\thesection}{1em}{}
\titleformat{\subsection}
  {\normalfont\large\bfseries\color{steelgrey}}
  {\thesubsection}{1em}{}

% ============ 超链接 ============
\usepackage{hyperref}
\hypersetup{
  colorlinks=true,
  linkcolor=navy,
  urlcolor=navy,
  citecolor=navy,
}

% ============ 其他工具 ============
\usepackage{enumitem}
\setlist{nosep, leftmargin=2em}
\usepackage{fontawesome5}
\usepackage{amssymb}
\usepackage{pifont}
\usepackage{tcolorbox}
\tcbuselibrary{skins,breakable}
\usepackage{caption}
\captionsetup{font=small, skip=4pt}
\usepackage{float}
\setlength{\textfloatsep}{8pt plus 2pt minus 2pt}
\setlength{\floatsep}{6pt plus 2pt minus 2pt}
\setlength{\intextsep}{6pt plus 2pt minus 2pt}
\setlength{\abovecaptionskip}{4pt}
\setlength{\belowcaptionskip}{2pt}

% ============ 自定义命令 ============
\newcommand{\rmark}{\textcolor{riskred}{\ding{108}}}
\newcommand{\amark}{\textcolor{riskamber}{\ding{108}}}
\newcommand{\gmark}{\textcolor{riskgreen}{\ding{108}}}
\newcommand{\starfive}{$\bigstar\bigstar\bigstar\bigstar\bigstar$}
\newcommand{\starfour}{$\bigstar\bigstar\bigstar\bigstar$}
\newcommand{\starthree}{$\bigstar\bigstar\bigstar$}

% 信息框
\newtcolorbox{keyinsight}[1][]{
  colback=lightgrey, colframe=navy,
  fonttitle=\bfseries, title={#1}, breakable
}
\newtcolorbox{riskbox}[1][]{
  colback=red!5, colframe=riskred,
  fonttitle=\bfseries, title={#1}, breakable
}
 after \documentclass[a4paper,11pt,openany,fontset=none]{ctexrep}

% ============ 页面布局 ============
\usepackage[top=2cm, bottom=2cm, left=2cm, right=2cm]{geometry}
\setlength{\parskip}{0.3em}
\setlength{\parindent}{2em}
\renewcommand{\baselinestretch}{1.15}

% ============ 字体（macOS） ============
\setCJKmainfont[BoldFont={Heiti SC}]{STSong}
\setCJKsansfont{Heiti SC}
\setCJKmonofont{Heiti SC}
\setmainfont{Times New Roman}

% ============ 颜色（IB Navy/Grey palette） ============
\usepackage{xcolor}
\definecolor{navy}{HTML}{003366}
\definecolor{steelgrey}{HTML}{4A5568}
\definecolor{lightgrey}{HTML}{F7FAFC}
\definecolor{riskred}{HTML}{C53030}
\definecolor{riskamber}{HTML}{D69E2E}
\definecolor{riskgreen}{HTML}{38A169}
\definecolor{nvgreen}{HTML}{76B900}

% ============ 表格 ============
\usepackage{booktabs}
\usepackage{longtable}
\usepackage{tabularx}
\usepackage{multirow}
\usepackage{array}
\newcolumntype{Y}{>{\centering\arraybackslash}X}
\newcolumntype{L}[1]{>{\raggedright\arraybackslash}p{#1}}
\newcolumntype{C}[1]{>{\centering\arraybackslash}p{#1}}
\newcolumntype{R}[1]{>{\raggedleft\arraybackslash}p{#1}}

% ============ 图形 ============
\usepackage{tikz}
\usetikzlibrary{shapes,arrows.meta,positioning,calc,backgrounds,fit}
\usepackage{pgfplots}
\usepgfplotslibrary{polar}
\pgfplotsset{compat=1.18}
\usepackage[edges]{forest}

% ============ 页眉页脚（报告标题和日期需在main.tex中覆盖） ============
\usepackage{fancyhdr}
\pagestyle{fancy}
\fancyhf{}
\fancyhead[L]{\small\color{steelgrey} \reporttitle}
\fancyhead[R]{\small\color{steelgrey} \reportdate}
\fancyfoot[C]{\small\color{steelgrey} \thepage}
\fancyfoot[R]{\tiny\color{steelgrey} 本报告不构成任何投资建议}
\renewcommand{\headrulewidth}{0.4pt}
\renewcommand{\footrulewidth}{0.2pt}

% ============ 章节格式 ============
\usepackage{titlesec}
\titleformat{\chapter}[display]
  {\normalfont\huge\bfseries\color{navy}}
  {\chaptertitlename\ \thechapter}{20pt}{\Huge}
\titleformat{\section}
  {\normalfont\Large\bfseries\color{navy}}
  {\thesection}{1em}{}
\titleformat{\subsection}
  {\normalfont\large\bfseries\color{steelgrey}}
  {\thesubsection}{1em}{}

% ============ 超链接 ============
\usepackage{hyperref}
\hypersetup{
  colorlinks=true,
  linkcolor=navy,
  urlcolor=navy,
  citecolor=navy,
}

% ============ 其他工具 ============
\usepackage{enumitem}
\setlist{nosep, leftmargin=2em}
\usepackage{fontawesome5}
\usepackage{amssymb}
\usepackage{pifont}
\usepackage{tcolorbox}
\tcbuselibrary{skins,breakable}
\usepackage{caption}
\captionsetup{font=small, skip=4pt}
\usepackage{float}
\setlength{\textfloatsep}{8pt plus 2pt minus 2pt}
\setlength{\floatsep}{6pt plus 2pt minus 2pt}
\setlength{\intextsep}{6pt plus 2pt minus 2pt}
\setlength{\abovecaptionskip}{4pt}
\setlength{\belowcaptionskip}{2pt}

% ============ 自定义命令 ============
\newcommand{\rmark}{\textcolor{riskred}{\ding{108}}}
\newcommand{\amark}{\textcolor{riskamber}{\ding{108}}}
\newcommand{\gmark}{\textcolor{riskgreen}{\ding{108}}}
\newcommand{\starfive}{$\bigstar\bigstar\bigstar\bigstar\bigstar$}
\newcommand{\starfour}{$\bigstar\bigstar\bigstar\bigstar$}
\newcommand{\starthree}{$\bigstar\bigstar\bigstar$}

% 信息框
\newtcolorbox{keyinsight}[1][]{
  colback=lightgrey, colframe=navy,
  fonttitle=\bfseries, title={#1}, breakable
}
\newtcolorbox{riskbox}[1][]{
  colback=red!5, colframe=riskred,
  fonttitle=\bfseries, title={#1}, breakable
}

\sloppy
\setlength{\emergencystretch}{3em}
\hypersetup{pdfauthor={\reportauthor},pdftitle={\reporttitle}}
\usepackage{amsmath}
\begin{document}
\astockcover
\begin{abstract}
本文回答一个严格限定的问题：从2026年7月22日到2026年12月31日，A股是否存在具备翻倍空间、且证据链足以进入投资研究核心池的标的。结论是：没有基准情景下的高置信度翻倍标的；雅化集团是唯一条件核心候选，恩捷股份因客户级量价与利用率证据不足只保留零权敏感性。原始报告中的42元与103元均具备数学上的翻倍空间，但目标期限未披露，不能改写成年末目标；其中只有雅化通过公司级估值资格门槛。报告将外部目标、AStock年末情景和证据阻断严格分开，并把缺失目标价、低基数、投资收益、周期峰值及现金流不匹配纳入降级规则。

\textbf{English abstract.} No validated name offers a high-confidence base-case double by 2026 year-end. 002497 is the sole conditional-core candidate, with a judgmental 20\% bull probability. 002812 retains only a zero-weight external sensitivity, with no house target or probability. The report separates reported evidence, broker forecasts, house inference, and timing risk.
\end{abstract}
\tableofcontents
\clearpage
\chapter{投资结论：有尾部空间，没有高置信度翻倍}

\begin{houseviewbox}[AStock House View]
截至2026年7月22日收盘，本报告没有找到“基准情景下、到2026年末高置信度翻倍”的A股标的。雅化集团是唯一进入\textbf{条件核心}研究层的候选；恩捷股份拥有可审计的103元券商目标，但关键量价、利用率和客户证据集中于单一券商假设，因此只列\textbf{零权敏感性备选}。中国船舶是基本面质量最好的零权备选；盛新锂能、江波龙与天华新能依次列为周期备选或观察。
\end{houseviewbox}

这个结论并不否认股价可能翻倍，而是严格区分三件事：第一，数学上的目标价是否超过现价一倍；第二，公司盈利是否足以支撑该目标；第三，市场能否在只剩五个月的窗口里完成重估。只满足第一项，不足以进入核心池。

\section{第一行动卡}

\begin{exhibitbox}[Exhibit 1：首要标的决策卡]
\small
\begin{tabularx}{\exhibitboxwidth}{L{2.2cm} R{1.3cm} L{2.2cm} L{1.7cm} L{1.4cm} X X}
\toprule
标的 & 现价 & 正式/基准状态 & 外部牛市参考 & 证据级别 & 下一验证 & 最强失效条件 \\
\midrule
雅化集团 & 16.79 & 约25--26元；条件核心 & 42元；期限未披露 & B+ & H2利润桥与OCF & 现金仍负或库存/套保主导 \\
恩捷股份 & 47.84 & 无正式目标；零权敏感性 & 103元；期限未披露 & C & 公司级量价、利用率与良率 & 关键参数继续仅见于单一券商 \\
\bottomrule
\end{tabularx}
\sourcenote{公司公告、2026年7月22日收盘价、原始券商报告与AStock估值治理。}
\qualitynote{外部目标期限均未披露；恩捷未通过公司级估值资格门槛，103元不进入正式目标。}
\end{exhibitbox}

\section{排名与行动图}

\begin{exhibitbox}[Exhibit 2：2026年末翻倍候选分层]
\small
\begin{tabularx}{\exhibitboxwidth}{C{0.8cm} L{1.7cm} L{2.4cm} R{1.35cm} R{1.35cm} R{1.35cm} X}
\toprule
排名 & 分层 & 标的 & 现价 & 基准/最终空间 & 牛市空间 & 研究结论 \\
\midrule
1 & 条件核心 & 002497 雅化集团 & 16.79 & 约+51\% & +150.1\% & 原始42元目标可审计；H2现金流与锂价决定能否升级 \\
2 & 备选 & 600150 中国船舶 & 33.02 & 无正式目标 & +108.7\% & 订单/利润质量较优，但无当前外部点目标 \\
3 & 零权敏感性 & 002812 恩捷股份 & 47.84 & 无正式目标 & +115.3\% & 103元目标可审计，但公司级量价证据阻断正式估值 \\
4 & 周期备选 & 002240 盛新锂能 & 27.65 & 无正式目标 & +113.5\% & 量价和自有矿弹性明确，现金流与锂价高度敏感 \\
5 & 观察 & 301308 江波龙 & 388.45 & 无正式目标 & +100.8\% & 峰值利润、负现金流与近期大幅回撤叠加 \\
6 & 观察 & 300390 天华新能 & 55.88 & 无正式目标 & +83.3\% & 当前牛市区间也未达到翻倍门槛 \\
\bottomrule
\end{tabularx}
\sourcenote{公司公告、2026年7月22日收盘价、原始券商报告、AStock情景模型。}
\qualitynote{概率与估值区间均为情景判断；缺显式点目标者零外部权重；显式目标也须先通过公司级估值资格门槛。}
\end{exhibitbox}

\section{两个“有空间”与一个“可进入核心”的区别}

雅化集团和恩捷股份的原始券商目标分别为42元和103元，相对现价隐含150.15\%和115.30\%上涨。但两份报告均未披露目标期限，因此它们不能被改写为“2026年末目标价”。本报告只把42元作为限权外部牛市锚；103元因恩捷未通过公司级估值资格门槛，仅保留为零权敏感性。

雅化集团能够进入条件核心，原因不是目标价最高，而是公司级证据更完整：2025年报披露了锂业务与民爆业务的收入、毛利率、锂产品销量、近13万吨产能、头部客户及重大合同正常履行。民爆业务提供稳定利润底仓，锂业务提供上行弹性。其主要缺口是H2实现价、库存收益、套保影响和经营现金流。

恩捷股份虽然有103元目标，但160亿平方米全年出货、单平利润、2027年利用率与有效产能等关键参数主要来自同一份券商模型，公司预告只确认了产销量增长、价格企稳和毛利改善。因此它只能做零权估值敏感性，暂不发布正式目标。

\begin{riskbox}[最容易误读的地方]
“具备翻倍空间”不等于“预期收益率100\%”。雅化集团的判断性翻倍概率约20\%，驱动情景与限权外部锚形成的审计值为25.32元，读者口径约25--26元。恩捷103元只是一份期限未披露的外部零权敏感性，不代表AStock给出了翻倍胜率或正式目标。
\end{riskbox}

\section{为什么不给出“必买名单”}

从当前到年末约五个月，价格翻倍对应约15\%的月复合涨幅。即使全年盈利预测正确，估值也必须快速重估。当前市场仍处于高成交、高波动和快速风格切换状态：7月21日创业板大涨7.05\%，次日又下跌3.23\%，7月22日全市场超过3,800只股票下跌。这样的市场有能力制造翻倍，也同样有能力制造错误的低估值幻觉。

因此，本报告的行动不是“追目标价”，而是把未来验证分为三层：盈利是否达到桥接、现金是否跟上利润、市场是否愿意支付更高倍数。三层缺一，牛市目标只能留在尾部情景。

\clearpage

\chapter{市场环境与五个月翻倍的算术}

\section{指数稳定掩盖了高波动个股环境}

2026年7月22日，上证指数报3,867.03点，上涨0.07\%；深证成指下跌1.42\%，创业板指下跌3.23\%。两市成交约2.65万亿元，但超过3,800只个股下跌。前一交易日创业板刚刚上涨7.05\%。这组数据说明市场并不缺流动性，却缺少稳定的风险偏好。

对翻倍筛选而言，这种环境具有双重意义。一方面，短期大幅回撤把部分盈利修复标的压到较低隐含倍数；另一方面，快速轮动降低了估值重估持续到年末的概率。低价本身不是安全边际，只有现金流与盈利证据可以把它转化为安全边际。

\begin{exhibitbox}[Exhibit 2：近十个交易日回撤审计]
\small
\begin{tabularx}{\exhibitboxwidth}{L{2.3cm} R{1.5cm} R{1.5cm} R{1.5cm} R{1.7cm} X}
\toprule
标的 & 7月9日收盘 & 7月22日收盘 & 区间涨跌 & 峰谷回撤 & 估值含义 \\
\midrule
中国船舶 & 35.695 & 33.02 & -7.49\% & -12.82\% & 相对防御，但仍需重估 \\
江波龙 & 620.00 & 388.45 & -37.35\% & -39.28\% & 高拥挤、高周期定价快速出清 \\
恩捷股份 & 59.33 & 47.84 & -19.37\% & -19.37\% & 市场显著折价2027恢复 \\
盛新锂能 & 35.55 & 27.65 & -22.22\% & -26.89\% & 锂价与风险偏好双重敏感 \\
天华新能 & 69.10 & 55.88 & -19.13\% & -22.82\% & 周期重估降温 \\
雅化集团 & 20.20 & 16.79 & -16.88\% & -19.36\% & 低倍数与盈利持续性怀疑并存 \\
\bottomrule
\end{tabularx}
\sourcenote{腾讯前复权日线；2026年7月9日至7月22日。中国船舶7月20日除息0.365元，价格比较并非总回报。}
\end{exhibitbox}

\section{翻倍所需的月度复合回报}

若从7月22日持有到12月31日，约五个多月内实现100\%回报，所需月复合涨幅约为：
\[
  r_m = 2^{1/5}-1 \approx 14.9\%.
\]
这不是普通的估值修复，而是“盈利兑现、预期上修、风险偏好改善”三者共振。只有一个环节发生，通常只能带来20\%--60\%的修复，而不是翻倍。

\section{当前价格隐含的怀疑}

雅化集团现价对应2026E EPS约6.38倍，翻倍到33.58元需要12.77倍；恩捷股份现价对应2027E EPS约9.27倍，翻倍到95.68元需要18.54倍。两者并非只要盈利不下修就能翻倍，还需要倍数显著上升。

中国船舶当前约对应2027E EPS 10.55倍，翻倍需要21.10倍。江波龙当前约对应2026E EPS 14.68倍，翻倍需要29.35倍，而券商预测2027年利润增长只有约2\%。这说明江波龙的翻倍逻辑更依赖估值而不是盈利，风险显著高于表面低PE。

\begin{keyinsight}[投资含义]
近期回撤让“目标价/现价”比值迅速上升，但分子没有同步变得更可信。筛选必须同时更新现价、盈利分母、目标期限和现金流，不能把旧目标价除以新低价后直接宣布机会。
\end{keyinsight}

\clearpage

\chapter{筛选方法：从数学翻倍到证据可投}

\section{五道门槛}

本报告先从2026H1盈利改善、低位估值、产业周期与已有模型中形成候选，再依次通过五道门槛：

\begin{enumerate}
  \item \textbf{盈利兑现：}2025年报、2026Q1和H1预告能否建立不机械年化的全年桥接；
  \item \textbf{估值空间：}基准与牛市情景是否分离，现价翻倍需要多少EPS与PE；
  \item \textbf{年末催化：}H2业绩、价格、产能、订单或整合是否能在12月31日前被市场观察；
  \item \textbf{现金与资产负债表：}利润是否转为经营现金流，库存和债务是否吞噬收益；
  \item \textbf{证据质量：}公司公告、原始券商报告、行情与推断是否清楚分层。
\end{enumerate}

评分权重为盈利兑现30\%、估值空间25\%、催化时点20\%、现金与资产负债表15\%、证据质量10\%。任何发布阻断项都可以让高分候选降为观察，而不是用总分掩盖证据缺口。

\begin{exhibitbox}[Exhibit 4：六标的可复核评分与硬门槛]
\scriptsize
\begin{tabularx}{\exhibitboxwidth}{L{1.9cm} R{1.0cm} R{1.0cm} R{1.0cm} R{1.0cm} R{1.0cm} R{1.0cm} L{1.1cm} X}
\toprule
标的 & 盈利30\% & 估值25\% & 催化20\% & 现金15\% & 证据10\% & 总分 & 证据级 & 硬门槛/最终层级 \\
\midrule
雅化集团 & 85 & 90 & 75 & 45 & 82 & 78.0 & B+ & OCF条件；条件核心 \\
中国船舶 & 78 & 55 & 68 & 75 & 80 & 70.0 & B+ & 缺点目标；质量备选 \\
恩捷股份 & 72 & 80 & 70 & 60 & 34 & 68.0 & C & 估值阻断；零权敏感性 \\
盛新锂能 & 70 & 70 & 65 & 30 & 50 & 61.0 & C+ & 周期/OCF；周期备选 \\
江波龙 & 75 & 55 & 55 & 20 & 38 & 54.1 & C & 峰值/OCF；观察 \\
天华新能 & 65 & 35 & 55 & 40 & 47.5 & 50.0 & C & 牛市不翻倍；观察 \\
\bottomrule
\end{tabularx}
\sourcenote{AStock评分：总分=各项分数乘页首权重；硬门槛独立于总分。}
\end{exhibitbox}

条件核心允许形成正式自有区间，但升级仍依赖已列验证项；质量备选与周期备选只表示研究优先级，不发布正式目标；观察表示驱动或现金证据不足；零权敏感性只展示外部或假设区间，权重为零。硬门槛覆盖估值资格、现金错配、周期峰值与翻倍阈值，不能被总分抵销。

\section{从16只候选到1只条件核心}

初始池覆盖中国船舶、江波龙、恩捷股份、盛新锂能、天华新能、雅化集团、赣锋锂业，以及若干低位控股平台、现金型公司、化工与AI硬件对照。最终六只进入深度经营与估值审查，其中只有雅化集团同时满足公司级经营证据、显式外部目标、当前价格与可复核情景。

\begin{exhibitbox}[Exhibit 3：筛选漏斗]
\centering
\begin{tikzpicture}[node distance=0.55cm, every node/.style={font=\small}]
\node[draw=navy,fill=lightgrey,rounded corners,minimum width=12.5cm,minimum height=0.8cm] (u) {16只初始候选：盈利预告、低位、周期修复与热门对照};
\node[below=of u,draw=navy,fill=lightgrey,rounded corners,minimum width=10.5cm,minimum height=0.8cm] (d) {6只深度审查：股本、现金流、经营驱动、券商原文、情景估值};
\node[below=of d,draw=riskamber,fill=yellow!8,rounded corners,minimum width=8.5cm,minimum height=0.8cm] (v) {2只有完整显式外部目标：雅化集团、恩捷股份};
\node[below=of v,draw=riskgreen,fill=green!6,rounded corners,minimum width=6.5cm,minimum height=0.8cm] (c) {1只条件核心：雅化集团};
\draw[-{Latex[length=3mm]},thick,navy] (u)--(d);
\draw[-{Latex[length=3mm]},thick,navy] (d)--(v);
\draw[-{Latex[length=3mm]},thick,navy] (v)--(c);
\end{tikzpicture}
\sourcenote{AStock筛选与证据门槛。}
\end{exhibitbox}

\section{外部目标的统一权重政策}

雅化集团42元目标以2026E EPS 2.63元乘16倍PE得出；恩捷股份103元目标以2027E EPS 5.16元乘20倍PE得出。两份原始报告完整、算术可复核，但均未披露目标期限，不能视为精确的2026年12月31日目标。

统一政策是：缺少显式点目标者，外部目标权重为零；显式目标但期限未披露者，只有在公司级经营驱动通过估值资格门槛后，才能作为上限10\%的牛市参考。雅化因此采用80\%驱动情景概率值、10\%现价市场锚、10\%原始外部目标；恩捷虽然有103元目标，但公司级量价、利用率与客户证据阻断估值资格，外部目标和自有正式目标权重均为零。

\section{不允许的估值捷径}

吉林敖东、辽宁成大等控股平台可能因投资收益或资产折价而出现高数学空间，但不能把投资收益当作高增长经营EPS；九安医疗的历史利润容易受一次性或波动项目影响；江波龙的2026利润可能处于存储涨价峰值，不能仅凭低PE给予增长倍数；AI硬件热门股也不能因为景气强就忽略估值拥挤。

\begin{disclosurebox}[证据分层]
公司公告用于确认已披露事实；原始券商PDF用于确认预测与目标；实时行情用于确认现价；AStock模型用于明确推断。任何客户、订单、ASP、产能利用率或单平利润若只来自券商估计，都必须标记为预测并折价。
\end{disclosurebox}

\clearpage

\chapter{条件核心：雅化集团（002497）}

\begin{dashboardbox}[结论卡]
\textbf{现价：}16.79元 \quad \textbf{总市值：}193.52亿元 \quad \textbf{正式区间：}约25--26元 \quad \textbf{牛市锚：}42元 \\
\textbf{基准/最终空间：}约+51\% \quad \textbf{牛市空间：}+150.1\% \quad \textbf{判断性翻倍概率：}20\% \\
\textbf{核心判断：}证据最完整的锂盐高弹性标的，民爆提供利润底仓；但FY2025和2026Q1经营现金流均为负，H2必须证明盈利不是库存与套保的短期组合。
\end{dashboardbox}

\section{为什么它是唯一条件核心}

雅化集团的优势不只是锂价弹性，而是双主业可以被分开验证。2025年锂业务收入48.24亿元，占总收入56.46\%，毛利率12.22\%；民爆业务收入33.16亿元，占38.82\%，整体毛利率37.46\%。锂业务决定上行斜率，民爆业务提供相对稳定的利润底盘，两者不能套用同一倍数。

公司年报披露当前锂盐综合产能近13万吨，2025年锂产品销量6.95万吨；长期合作客户包括Tesla、LGES、SK ON、宁德时代等，重大合同处于正常履行状态。头部客户收入占锂业务90\%以上。这些信息比仅有行业需求或客户认证的标的更接近公司级收入证据。

民爆业务同样有真实经营壁垒：牌照、区域渠道、项目服务和安全运营构成进入门槛。2025年工业炸药许可产能26.25万吨、利用率83.74\%；2026年4月许可产能增至33万吨。新增许可不等于新增利润，仍需观察利用率与回款。

\begin{exhibitbox}[Exhibit 4：双主业经营桥]
\small
\begin{tabularx}{\exhibitboxwidth}{L{2.3cm} L{3.2cm} L{3.2cm} X}
\toprule
分部 & 已披露经营证据 & 2026年关键预测/推断 & 估值处理 \\
\midrule
锂盐与资源 & 近13万吨产能；2025销量6.95万吨；头部客户与重大合同正常履行 & 券商预计全年约12万吨出货、自供比例约25\%；Q2吨利2.6--3.4万元 & 周期PE；剔除库存与套保一次影响 \\
民爆 & 2025收入33.16亿元、毛利率37.46\%；炸药利用率83.74\% & 券商预计利润约5.5--6亿元；新增许可产能需验证利用率 & 稳定分部信用，不给锂盐高弹性倍数 \\
运输与其他 & 2025收入4.03亿元 & 不作为核心增长来源 & 不单独给溢价 \\
\bottomrule
\end{tabularx}
\sourcenote{公司2025年报、2026H1业绩预告；东吴证券2026年7月7日原始报告。}
\end{exhibitbox}

\section{盈利从预告到全年还差什么}

公司2025年收入85.43亿元、归母净利润6.32亿元；2026Q1收入28.30亿元、归母3.39亿元、扣非3.57亿元。H1预告归母净利润11--13亿元、扣非11.25--13.15亿元。与券商2026全年30.32亿元归母净利润相比，H2仍需贡献17.32--19.32亿元。

这个H2要求不是简单的季节性外推，而是量、价、自供矿、库存与套保共同作用。原始报告预计2026年锂盐出货约12万吨、权益资源自供比例约25\%，并认为套保拖累在下半年减弱。公司公告只确认H1销量与销售均价同步增长，没有披露吨位、实现ASP与单位成本；因此吨利和出货必须继续标记为券商预测。

最关键的反证来自现金流。FY2025经营现金流为-5.70亿元，2026Q1为-4.31亿元。券商还预计2026年存货可能从16.96亿元增至51.13亿元。若利润增长主要来自库存价差，而现金继续被营运资金占用，则低PE并不代表可分配价值。

\begin{exhibitbox}[Exhibit 5：盈利—现金验证表]
\small
\begin{tabularx}{\exhibitboxwidth}{L{2.5cm} R{2cm} R{2cm} R{2cm} X}
\toprule
指标 & FY2025 & 2026Q1 & 2026H1预告 & 研究含义 \\
\midrule
营业收入 & 85.43亿元 & 28.30亿元 & 未披露 & 不从利润倒推收入 \\
归母净利润 & 6.32亿元 & 3.39亿元 & 11--13亿元 & H2仍需17.3--19.3亿元达到券商全年预测 \\
扣非净利润 & 6.94亿元 & 3.57亿元 & 11.25--13.15亿元 & 扣非改善较纯投资收益更可信 \\
经营现金流 & -5.70亿元 & -4.31亿元 & 未披露 & 核心阻断项，要求H2明显转正 \\
\bottomrule
\end{tabularx}
\sourcenote{公司2025年报、2026Q1报告、2026H1预告。H1预告未经审计。}
\end{exhibitbox}

\section{估值：为什么42元只属于牛市情景}

原始券商模型用2026E EPS 2.63元乘16倍PE，得到42.08元并给出42元目标。现价只对应6.38倍。若到年末翻倍至33.58元，需要12.77倍2026E PE，低于券商16倍但约为当前倍数的两倍。

本报告不让三档情景共用同一个EPS。悲观/基准/牛市分别假设锂盐出货9/12/13万吨、归母净利润13.1/30.35/38.5亿元、EPS 1.14/2.63/3.34元，再给予8/10/12.6倍PE，对应9.12/26.30/42.08元。概率分别为30\%/50\%/20\%，经营驱动概率值24.30元；再与现价和期限未披露的42元外部锚按80\%/10\%/10\%组合，审计值为25.32元，读者口径约25--26元。该方法明确表达：盈利兑现可以带来较大修复，但翻倍还需要更强盈利与倍数同时出现。

\begin{riskbox}[降级条件]
若Q3经营现金流继续明显为负、利润主要来自库存收益或套保、锂盐出货与自供比例未兑现，或民爆回款与利润恶化，雅化集团应从条件核心降为周期观察。价格继续下跌不是自动升级理由。
\end{riskbox}

\section{年末前的必要证据}

三季报需要看到锂盐出货接近3万吨、剔除库存与套保后吨利至少约2万元、经营现金流转正或显著改善；全年需要接近12万吨出货、归母利润接近30亿元、现金流转正、民爆利润约5.5--6亿元。吨位与吨利来自券商估计，现金阈值属于AStock验证要求。

只有当这些证据共同出现，市场把2026E EPS从6倍重估到13倍以上才具备基本面基础。若只出现锂价上涨而现金与自供比例不改善，42元仍只是周期尾部。

\clearpage

\chapter{备选：恩捷股份（002812）}

\begin{dashboardbox}[结论卡]
\textbf{现价：}47.84元 \quad \textbf{总市值：}469.46亿元 \quad \textbf{正式目标：}不发布 \quad \textbf{外部零权敏感性：}103元 \\
\textbf{基准/最终空间：}不发布 \quad \textbf{外部目标相对现价：}+115.3\% \quad \textbf{AStock翻倍概率：}不估计 \\
\textbf{核心判断：}盈利修复真实、现金流为正，且原始103元目标可审计；但目标建立在2027E利润跃升和单一券商的量价/利用率假设上，未通过正式估值资格门槛。
\end{dashboardbox}

\section{已经确认的修复}

恩捷股份2025年收入136.33亿元、归母净利润1.43亿元；2026Q1收入39.08亿元、归母2.60亿元、扣非2.61亿元，毛利率由2025年的18.77\%改善至28.09\%。H1预告归母7.36--9.00亿元、扣非7.62--9.30亿元，说明扭亏与毛利修复已进入报表。

现金质量相对锂盐和存储标的更好。FY2025经营现金流11.44亿元，2026Q1为2.07亿元，均为正。不过，隔膜业务资本开支高，新产线投放与利用率爬坡仍可能重新占用现金，不能仅以当前正现金流推断2027年盈利。

公司官方预告确认下游需求旺盛、产销量增长、价格较上年低点企稳回升、综合毛利率显著改善。官方没有披露H1出货量、客户涨价、实际单平利润、有效产能、利用率与良率，这些正是103元目标的关键变量。

\section{券商模型的强项与集中风险}

东吴证券预计2026年出货超过160亿平方米，2027年有效产能约220亿平方米，并假设湿法隔膜利用率从2026年的90\%提升至2027年的97\%；其2026/2027/2028归母净利润预测为22.79/50.68/62.01亿元，对应EPS 2.32/5.16/6.31元。目标价103元来自2027E EPS 5.16元乘20倍PE。

该模型的优点是经营桥清楚：出货平方米、单平价格、单位利润、利用率与折旧摊薄能够连接到EPS。缺点是关键参数高度同源。客户名称、涨价比例、订单、有效产能与新线良率没有公司级量化证据。一个完整的券商模型并不等于多个独立证据。

\begin{exhibitbox}[Exhibit 6：隔膜恢复桥的证据层级]
\small
\begin{tabularx}{\exhibitboxwidth}{L{3cm} L{3.4cm} L{3.4cm} X}
\toprule
变量 & 公司已披露 & 券商预测 & 估值处理 \\
\midrule
需求/价格 & 需求旺盛；价格企稳回升 & 部分大客户、海外与3C涨价仍在落地 & 方向给信用，幅度折价 \\
出货量 & 产销量增长，无吨位 & H1 76--77亿平；全年160亿平以上 & 不作为已实现事实 \\
单平利润 & 未披露 & Q2约0.15元；2027约0.20--0.25元 & Bull核心参数，需Q3验证 \\
利用率/产能 & 未披露同口径有效产能 & 2026/2027利用率90\%/97\%；2027有效产能220亿平 & 高度敏感，未验证前限制倍数 \\
现金流 & FY2025与Q1均为正 & 扩产资本开支仍较高 & 需要OCF覆盖扩产 \\
\bottomrule
\end{tabularx}
\sourcenote{公司公告；东吴证券2026年7月10日原始报告。}
\end{exhibitbox}

\section{估值：2027盈利必须被提前资本化}

当前价格对应2027E EPS约9.27倍。翻倍至95.68元需要18.54倍，接近券商20倍假设。也就是说，市场不仅要相信利润从2026年的22.79亿元增至2027年的50.68亿元，还要在2026年末前给予接近完整恢复周期的倍数。

东吴证券2027E EPS 5.16元乘6倍、13倍和20倍，可机械得到30.96元、67.08元和103.20元。三档全部依赖同一券商盈利分母，且公司尚未披露足以验证出货、单平利润、利用率、良率与客户涨价的参数，因此本报告不为其赋概率、不混合现价，也不发布正式目标。该区间只回答“若卖方分母兑现，估值可能落在哪里”，不构成AStock预期收益。

\section{为什么排名落后于中国船舶}

恩捷拥有显式目标，而中国船舶没有；但研究排名不等于目标价排名。中国船舶的订单与交付能见度、Q1现金流和利润质量更强，经营桥中的公司级缺口更少。恩捷的客户涨价、利用率、良率与单平利润尚未公开量化，因此在核心资格上必须让位于证据质量。

\begin{riskbox}[降级条件]
若Q3单平利润低于约0.12元、客户涨价继续延迟、新线良率或利用率不及预期、经营现金流被应收与资本开支吞噬，103元牛市情景应删除。只有出货约40亿平以上、单平利润约0.15元、客户涨价与现金流共同验证，才可重审核心资格。
\end{riskbox}

\clearpage

\chapter{备选与观察：为什么暂不升级}

\section{中国船舶：质量最好的零权备选}

中国船舶2025年收入1,519.78亿元、归母78.48亿元；2026Q1收入433.12亿元、归母48.32亿元、扣非47.67亿元、经营现金流81.52亿元。H1预告归母92--110亿元、扣非90--108亿元。利润、扣非与现金三者同时改善，是六只深度标的中最扎实的报表组合。

原始券商报告预计2026--2028年归母净利润181.98/235.38/299.43亿元，EPS 2.42/3.13/3.98元，并判断交付排期延伸至2028--2030年。订单能见度和高价订单进入交付构成核心逻辑。需要注意的是，2025年资产整合和同一控制下合并会扭曲同比增速，不能直接把并表增长套入高成长PE。

现价33.02元约对应2027E EPS 10.55倍，翻倍需要21.10倍。AStock牛市情景68.9元可以超过翻倍门槛，但当前收集的五份原始报告均没有显式目标价。不能从报告列示的当前PE反推券商目标，因此该情景保持零外部权重。

\begin{keyinsight}[升级条件]
若三季报毛利率维持约16\%以上、经营现金流持续为正、高价与中高端船交付继续，同时取得当前、显式、可审计的外部目标，中国船舶可以重新竞争条件核心。它的优势是质量，缺点是五个月内将倍数翻倍的难度。
\end{keyinsight}

\section{盛新锂能：商品周期尾部}

盛新锂能2025年归母亏损8.88亿元，2026Q1归母转为4.64亿元，H1预告归母10--12亿元、扣非13--15亿元。反转由锂盐价格、出货、自有矿与印尼工厂共同驱动，券商预测2026E EPS 2.45元、2027E EPS 3.28元。

现价27.65元翻倍至55.30元，需要22.57倍2026E EPS；若用2027E EPS，仍需16.86倍。AStock牛市情景59.0元能够翻倍，但原始报告没有显式目标，且2026Q1经营现金流为-6.67亿元。量价与自供比例必须和现金共同验证，单一锂价上涨不足以升级。

需要观察Q3出货是否接近3万吨、综合吨利是否接近2万元、印尼利用率是否继续提升、经营现金流是否转正。若全年出货接近12万吨、自有矿贡献和现金均可核验，才有资格从周期备选上调。

\section{江波龙：利润很强，但分母可能在峰值}

江波龙2026Q1收入99.09亿元、归母38.62亿元、扣非39.43亿元，毛利率55.53\%；H1预告收入220--250亿元、归母92--110亿元。利润增长极强，但Q1经营现金流为-28.75亿元，FY2025经营现金流也为-12.01亿元。

国信证券预测2026E净利润110.97亿元、EPS 26.47元，而2027E净利润仅增长约1.9\%。另一份原始报告对2026E净利润预测高达199.69亿元，两份报告在类似收入预测下出现88.72亿元利润差，说明晶圆价格、库存成本层与毛利假设主导模型。如此大的分歧不适合建立精确目标。

股价从7月9日620元跌至7月22日388.45元，十个交易日下跌37.35\%。现价翻倍需要约29.35倍2026E EPS。若2027利润不增长，市场必须在峰值分母上给出接近30倍PE，难度与回撤风险都很高。

\begin{riskbox}[核心阻断]
公司尚未披露bit出货、产品ASP、库存收益拆分、封测利用率、下游订单转量与客户收入。平台认证与LTA/MOU能够说明技术和供给关系，却不能证明可持续单位经济。只有单季经营现金流转正、库存周转改善、H2毛利稳定并出现2027年两位数利润增长，才应重建估值。
\end{riskbox}

\section{天华新能：盈利强，但当前价格已使翻倍门槛更高}

天华新能H1预告归母22--24亿元、扣非21.62--23.62亿元，Q1毛利率46.99\%，盈利恢复显著。券商预测2026E EPS约5.70元，AStock情景为56.4/79.8/102.4元。以现价55.88元计算，牛市空间约83.3\%，仍未达到翻倍所需的111.76元。

公司与宁德时代的股权关系可说明产业协同，不能自动推导采购订单和价格溢价；未投产矿山与在建二期产能也不进入2026现金流。现阶段它是有效的锂周期观察标的，但不是本题要求的翻倍候选。

\begin{exhibitbox}[Exhibit 7：零权备选的阻断项]
\small
\begin{tabularx}{\exhibitboxwidth}{L{2.4cm} R{1.5cm} R{1.5cm} X L{3.2cm}}
\toprule
标的 & 牛市空间 & Q1 OCF & 主要优势 & 发布阻断 \\
\midrule
中国船舶 & +108.7\% & +81.52亿元 & 订单、交付、现金质量 & 无显式外部目标；并表归一化 \\
盛新锂能 & +113.5\% & -6.67亿元 & 锂价、自有矿、印尼增量 & 无显式目标；商品与现金风险 \\
江波龙 & +100.8\% & -28.75亿元 & 存储价格与产品升级 & 峰值分母、预测分歧、现金缺口 \\
天华新能 & +83.3\% & -1.83亿元 & 锂盐量价与资源布局 & 牛市区间本身不翻倍 \\
\bottomrule
\end{tabularx}
\sourcenote{公司报告、H1预告、原始券商报告、AStock情景模型。}
\end{exhibitbox}

\clearpage

\chapter{估值：把尾部目标还原为所需条件}

\section{正式估值只覆盖一个合格标的}

正式估值表只纳入雅化集团，因为它同时具备当前价格、官方财务、公司级经营驱动、完整预测、显式外部目标、估值方法与原始PDF。恩捷虽然有显式目标，但公司级量价与客户证据触发估值阻断；中国船舶、江波龙、盛新锂能与天华新能则缺少点目标或未达到翻倍门槛，均不进入正式目标权重。

\begin{exhibitbox}[Exhibit 8：三档价值与正式目标]
\small
\begin{tabularx}{\exhibitboxwidth}{L{2.4cm} R{1.4cm} R{1.4cm} R{1.4cm} R{1.4cm} R{1.7cm} R{1.7cm} X}
\toprule
标的 & 现价 & 悲观 & 基准 & 牛市 & 牛市空间 & 正式目标 & 正式空间 \\
\midrule
雅化集团 & 16.79 & 9.12 & 26.30 & 42.08 & +150.6\% & 25.32 & +50.8\% \\
\bottomrule
\end{tabularx}
\sourcenote{公司股本、2026年7月22日收盘价、原始券商预测、AStock估值模型。}
\end{exhibitbox}

\section{雅化集团估值桥}

雅化集团不再让三档情景共用同一EPS。AStock把民爆利润底仓与锂业务分开，并明确剔除库存收益与套保损益：悲观/基准/牛市分别假设锂盐出货9/12/13万吨、锂业务正常化净贡献0.9/2.05/2.50万元每吨、民爆归母贡献5.0/5.75/6.0亿元，对应归母净利润13.1/30.35/38.5亿元及EPS 1.14/2.63/3.34元。再分别给予8/10/12.6倍PE，得到9.12、26.30、42.08元。

\begin{exhibitbox}[Exhibit 8A：雅化经营驱动至EPS桥]
\scriptsize
\begin{tabularx}{\exhibitboxwidth}{L{1.5cm} R{1.2cm} R{1.4cm} R{1.3cm} R{1.3cm} R{1.3cm} R{1.4cm} R{1.2cm} R{1.2cm} R{1.2cm}}
\toprule
情景 & 锂盐万吨 & 锂价代理万元/吨 & 收入亿元 & 毛利率 & 毛利亿元 & 费税等亿元 & 归母亿元 & EPS元 & PE/价值 \\
\midrule
悲观 & 9 & 12.0 & 120.0 & 20.0\% & 24.0 & 10.9 & 13.1 & 1.14 & 8x/9.12 \\
基准 & 12 & 16.5 & 170.52 & 27.04\% & 46.11 & 15.76 & 30.35 & 2.63 & 10x/26.30 \\
牛市 & 13 & 18.0 & 195.0 & 28.5\% & 55.58 & 17.08 & 38.5 & 3.34 & 12.6x/42.08 \\
\bottomrule
\end{tabularx}
\sourcenote{基准收入、毛利率、出货与民爆利润参考东吴证券；悲观/牛市为AStock情景假设。费税等为毛利至归母的合并桥。}
\qualitynote{吨贡献、锂价代理及牛熊参数不是公司指引；库存收益与套保损益均不计入正常化贡献。}
\end{exhibitbox}

三档概率为30\%/50\%/20\%，经营驱动概率值为24.30元，而不是由同一EPS只改PE得到。审计值计算为：
\[
0.8\times24.302+0.1\times16.79+0.1\times42.00=25.32\text{元}.
\]
正式空间为50.80\%。对基准EPS而言，翻倍门槛33.58元等于12.77倍；这个倍数并不极端，但从当前6.38倍提升至12.77倍仍需要现金流与周期持续性共同证明。读者应使用约25--26元的区间，而非把25.32元视作精确价格。

\section{恩捷股份：零权敏感性，不是正式估值桥}

东吴证券以2027E正常化EPS 5.16元乘20倍PE给出103元目标；6/13/20倍可机械得到30.96/67.08/103.20元。由于出货、客户涨价、实际单平利润、有效利用率和新线良率缺少公司级量化证据，这一组数值全部属于同源卖方敏感性。本报告不为其赋情景概率，不与现价混合，也不发布正式目标。翻倍门槛95.68元等于18.54倍该卖方EPS，说明市场几乎必须完整接受20倍恢复估值。

\section{为什么不用PEG或简单低PE}

2026年利润增速普遍建立在低基数或周期反转上：雅化从低锂价和套保拖累中恢复，恩捷从隔膜价格与利用率低谷恢复，江波龙受存储涨价与库存成本层影响，盛新与天华同样依赖锂价。PEG会把低基数增长误判为结构增长；简单PE则可能把峰值利润误判为永久分母。

本报告优先使用单位经济与现金流校准后的PE。雅化的民爆分部只获得稳定信用，不套锂盐倍数；恩捷的2027EPS只有在出货、单平利润、利用率与现金共同形成公司证据后才有资格进入正式模型；控股平台投资收益和一次性利润不进入成长PE。

\section{市值与股本审计}

雅化集团总股本11.5256亿股，现价对应总市值193.52亿元。恩捷股份在2026年7月21日完成81.34万股限制性股票注销，最新总股本9.8132亿股，现价对应总市值469.46亿元。若沿用Q1旧股本，会轻微高估市值并造成EPS口径漂移。

\begin{sourcequalitybox}[估值治理]
缺少显式点目标者零外部权重；显式目标但期限未披露者，只有通过公司级估值资格门槛才可获得上限10\%的牛市参考权重。雅化符合该条件，恩捷不符合。概率是审慎表达工具，不是统计胜率；正式输出应读作区间，而非精确点位。
\end{sourcequalitybox}

\clearpage

\chapter{催化日历：从H1落地到年末验证}

\section{催化必须在时间窗内发生}

本题的截止日是2026年12月31日，因此长期产能、2028年矿山投产或2030年订单能见度只能解释资产质量，不能替代年末催化。真正有用的催化集中在H1正式报告、三季报、产品价格与出货、经营现金流以及下一年度预期形成。

\begin{exhibitbox}[Exhibit 9：2026年下半年验证时间轴]
\centering
\begin{tikzpicture}[x=1cm,y=1cm,font=\small]
\draw[very thick,navy] (0,0)--(14,0);
\foreach \x/\lab in {0/7月下旬,3/8月,6/9月,9/10月,12/11月,14/12月末}{
  \draw[navy] (\x,0.18)--(\x,-0.18);
  \node[below=0.25cm] at (\x,0) {\lab};
}
\node[above=0.45cm,align=center,text width=2.8cm] at (2,0) {H1正式报告\\预告质量复核};
\node[above=1.35cm,align=center,text width=3cm] at (6,0) {锂盐/隔膜价格、出货\\新线与库存证据};
\node[above=0.45cm,align=center,text width=3cm] at (9,0) {三季报\\利润与OCF验证};
\node[above=1.35cm,align=center,text width=3cm] at (12.5,0) {2027预期形成\\年末估值选择};
\end{tikzpicture}
\sourcenote{AStock研究框架；具体定期报告日期以交易所安排为准。}
\end{exhibitbox}

\section{雅化集团：最短验证路径}

H1正式报告首先需要解释利润来源：销量、实现价、自供矿、库存与套保各自贡献多少。三季报需要看到锂盐出货接近3万吨、剔除库存与套保后的吨利至少约2万元、经营现金流转正或大幅改善。到年末，市场还需要相信全年约12万吨出货、约30亿元利润与民爆5.5--6亿元稳定利润能够持续。

如果利润兑现但现金仍差，基准目标可以成立，翻倍概率不应上调；如果利润、现金、自供比例和合同履行同时改善，12.77倍年末PE才有合理基础。

\section{恩捷股份：验证2027而非只验证H1}

H1正式报告只能确认扭亏。决定103元目标的是2027年盈利跃升，因此Q3必须出现客户涨价、出货约40亿平方米、单平利润约0.15元、新线良率与经营现金流等公司级证据。若这些参数仍只存在于券商模型，市场可能继续把2027E EPS按低倍数折现。

\section{中国船舶、盛新锂能与江波龙}

中国船舶需要高价船交付、毛利率约16\%以上和正经营现金流；盛新锂能需要约3万吨季度出货、约2万元综合吨利、印尼利用率与正现金流；江波龙需要单季现金转正、库存周转改善、毛利稳定和2027年增长重新加速。

\begin{exhibitbox}[Exhibit 10：升级与删除条件]
\small
\begin{tabularx}{\exhibitboxwidth}{L{2.0cm} X L{2.4cm} X}
\toprule
标的 & 升级证据 & 阈值性质 & 删除/降级证据 \\
\midrule
雅化集团 & H2利润桥、OCF转正、自供与民爆稳定 & AStock阈值；部分输入源于券商 & 现金持续为负、库存收益主导、锂价反转 \\
恩捷股份 & 客户涨价、单平利润、利用率/良率与OCF & 券商估计转为公司证据 & 单平低于约0.12元、涨价延迟、扩产吞噬现金 \\
中国船舶 & 高价交付、毛利、OCF及显式外部目标 & AStock阈值；经营输入源于券商 & 交付延迟、毛利回落、整合未兑现 \\
盛新锂能 & 出货、吨利、自有矿、印尼与OCF共同验证 & 券商估计/AStock阈值 & 锂价回落、资源贡献不足、短债/现金恶化 \\
江波龙 & 现金转正、库存改善、2027利润增长两位数 & AStock阈值；预测源于券商 & 毛利回落、库存继续上升、现金仍显著为负 \\
天华新能 & 资源自供与OCF提升且价格回到安全区间 & AStock阈值 & 牛市价值仍低于翻倍门槛 \\
\bottomrule
\end{tabularx}
\qualitynote{券商派生阈值不是公司指引或已实现事实；只有定期报告、公告或可审计经营数据能关闭验证缺口。}
\end{exhibitbox}

\clearpage

\chapter{反方与风险：翻倍情景如何失败}

\section{最强反方观点}

最强反方不是“这些公司没有增长”，而是市场可能正确地拒绝为周期利润支付高倍数。雅化与盛新的锂盐利润受商品价格、库存和套保影响；恩捷的2027利润依赖价格、利用率与客户验证；江波龙的存储利润可能处于涨价与低价库存共振的峰值；中国船舶的高价订单兑现虽然更长，但仍属于周期交付。

若分母回落，低PE会迅速变成高PE；若分母兑现但倍数不升，股价也无法在五个月内翻倍。年末目标必须同时承受盈利风险与估值风险。

\begin{exhibitbox}[Exhibit 11：概率—影响风险矩阵]
\small
\begin{tabularx}{\exhibitboxwidth}{L{3cm} C{1.4cm} C{1.4cm} L{2.8cm} X}
\toprule
风险 & 概率 & 影响 & 暴露标的 & 触发与动作 \\
\midrule
锂价反转 & 中高 & 高 & 雅化、盛新、天华 & 立即下修EPS与倍数，删除纯价格驱动Bull \\
利润不转现金 & 高 & 高 & 雅化、江波龙、盛新、天华 & H2 OCF仍负则降级或维持零权 \\
隔膜恢复停滞 & 中 & 高 & 恩捷 & 价格/利用率/良率未达标则下修2027E \\
存储峰值回落 & 高 & 高 & 江波龙 & 毛利与库存转差则拒绝高PE \\
船舶交付或毛利延迟 & 中 & 中高 & 中国船舶 & 下修2027EPS与整合信用 \\
风险偏好压缩 & 高 & 高 & 全部 & 延长兑现期，不沿用年末牛市目标 \\
目标期限未披露 & 高 & 中 & 雅化、恩捷 & 仅雅化通过资格后限权10\%；恩捷零权 \\
\bottomrule
\end{tabularx}
\end{exhibitbox}

\section{现金流是共同薄弱环节}

江波龙2026Q1经营现金流-28.75亿元、盛新锂能-6.67亿元、雅化集团-4.31亿元、天华新能-1.83亿元。雅化与江波龙FY2025经营现金流也为负。盈利上行而现金下行，可能来自库存、预付款、应收、采购保障或扩产；在周期反转时，这些项目会放大风险。

恩捷与中国船舶的现金表现相对更好。恩捷FY2025与Q1均为正，但扩产资本开支仍需覆盖；中国船舶Q1经营现金流81.52亿元，但船舶建造现金流受里程碑回款影响，单季数据不能永久外推。

\section{证据风险与卖方选择偏差}

本报告归档的14份原始券商报告全部为正面评级，没有中性或卖出报告。这不是市场真实评级分布，而是公开可获得报告的选择偏差。雅化与恩捷各只有一份完整点目标；恩捷的关键参数尤其集中于单一来源。

因此，报告不以“券商都看多”作为结论，而是把原始报告拆成可审计的预测、方法、目标与风险，再与公司公告和现金流交叉。目标价是一个假设包，而不是外部真理。

\section{期限风险}

东吴证券两份报告未披露目标价期限；其公司评级定义不能替代具体目标期限披露。42元和103元是否针对12月31日无法确认。将目标价直接贴上“年末”标签，会夸大确定性。

\section{结论失败的明确信号}

若雅化H2经营现金流不转正或利润主要来自库存，唯一条件核心资格取消；若恩捷无法提供客户涨价、单平利润与利用率证据，103元继续保持零权敏感性；若中国船舶不能保持毛利与现金，质量备选下调；若江波龙H2利润与现金继续背离，任何翻倍区间删除。

\begin{riskbox}[风险纪律]
本报告不设置机械止损或交易指令。研究上的“止损”是证据失效：当盈利桥、现金桥或估值期限不再成立，应更新结论，而不是等待股价替研究者证明原叙事正确。
\end{riskbox}

\section{组合压力测试}

压力测试不是预测，而是检验结论能否承受同时冲击。若锂价从当前中枢下跌25\%，雅化与盛新的盈利分母和估值倍数可能同步下修；若隔膜涨价延后两个季度，恩捷的2027E EPS与市场前置定价都需下调；若全市场高弹性板块PE再压缩30\%，即使盈利不变，所有年末翻倍情景都将显著后移。

\begin{exhibitbox}[Exhibit 12：三类压力与结论变化]
\small
\begin{tabularx}{\exhibitboxwidth}{L{3cm} L{3.6cm} X X}
\toprule
压力情景 & 首要受影响变量 & 结论变化 & 恢复条件 \\
\midrule
锂价中枢下跌25\% & 雅化/盛新/天华吨利、库存与PE & 雅化取消条件核心；锂链Bull删除 & 自供降本与现金流抵消价格冲击 \\
隔膜恢复延后两季度 & 恩捷单平利润、利用率、2027E EPS & 103元零权敏感性继续不进入正式模型 & 客户涨价与新线良率形成公司证据 \\
高弹性PE压缩30\% & 全部标的市场倍数 & 年末翻倍概率接近个位数 & 盈利上修足以抵消倍数收缩 \\
存储毛利回落15个百分点 & 江波龙利润与库存价值 & 观察区间下修，禁止使用峰值EPS & 现金转正、库存周转与产品结构改善 \\
\bottomrule
\end{tabularx}
\end{exhibitbox}

\clearpage

\chapter{证据质量与可审计边界}

\section{行情与股本}

16只初始候选的2026年7月22日收盘价均经AStock/Sina与腾讯行情双源匹配。六只深度标的的总股本均与公司公告或季度报告对齐。恩捷股份使用限制性股票注销后的9.8132亿股，江波龙使用限制性股票归属后的4.2306亿股。

行情质量足以确认现价与市值，但不能识别卖方身份、融资渠道或下跌原因。近十日回撤只用于刻画估值压缩，不自动解释为基本面恶化。

\section{财务与预告}

六只深度标的的2025年报与2026Q1收入、归母、扣非与经营现金流均与官方报告核对。H1预告直接来自公司公告，但均未经审计；除江波龙披露H1收入区间外，其余公司没有在预告中披露收入。锂盐与隔膜公司也未披露H1吨位、ASP、单位成本或单平利润。

\begin{exhibitbox}[Exhibit 13：官方财务锚]
\small
\begin{tabularx}{\exhibitboxwidth}{L{2.4cm} R{1.7cm} R{1.7cm} R{1.7cm} R{1.7cm} X}
\toprule
标的 & 2025收入 & Q1归母 & Q1 OCF & H1归母预告 & 关键质量结论 \\
\midrule
中国船舶 & 1,519.78亿 & 48.32亿 & +81.52亿 & 92--110亿 & 并表重述但现金强 \\
江波龙 & 227.66亿 & 38.62亿 & -28.75亿 & 92--110亿 & 利润/现金严重背离 \\
恩捷股份 & 136.33亿 & 2.60亿 & +2.07亿 & 7.36--9.00亿 & 扭亏已确认，单位经济未披露 \\
盛新锂能 & 50.64亿 & 4.64亿 & -6.67亿 & 10--12亿 & 周期反转，现金待验证 \\
天华新能 & 75.49亿 & 9.69亿 & -1.83亿 & 22--24亿 & 利润强，现金与资源自供待验证 \\
雅化集团 & 85.43亿 & 3.39亿 & -4.31亿 & 11--13亿 & 公司证据完整，现金仍是阻断 \\
\bottomrule
\end{tabularx}
\sourcenote{公司2025年报、2026Q1报告、2026H1预告。H1预告未经审计。}
\end{exhibitbox}

\section{券商证据}

本次归档14份原始PDF：中国船舶5份、江波龙2份、恩捷股份1份、盛新锂能4份、天华新能1份、雅化集团1份。只有雅化和恩捷的报告同时披露目标、预测与估值方法。其他报告的盈利预测仍可比较，但目标价权重为零。

雅化目标42元对应2026E EPS 2.63元乘16倍PE；恩捷目标103元对应2027E EPS 5.16元乘20倍PE。两份报告的目标期限均未披露，因而不支持“卖方预计年末翻倍”的表述。

\section{事实、预测与推断}

\begin{exhibitbox}[Exhibit 14：证据金字塔]
\centering
\begin{tikzpicture}[font=\small]
\fill[navy!85] (-1.7,3.0)--(1.7,3.0)--(0,4.1)--cycle;
\node[text=white,align=center] at (0,3.45) {公司公告\\已披露事实};
\fill[accentblue!75] (-3.0,1.9)--(3.0,1.9)--(1.7,3.0)--(-1.7,3.0)--cycle;
\node[text=white,align=center] at (0,2.45) {原始券商PDF\\预测、方法、目标};
\fill[steelgrey!65] (-4.2,0.8)--(4.2,0.8)--(3.0,1.9)--(-3.0,1.9)--cycle;
\node[text=white,align=center] at (0,1.35) {实时行情与结构化数据\\现价、市值、历史交易};
\fill[mutedgrey!35] (-5.2,-0.4)--(5.2,-0.4)--(4.2,0.8)--(-4.2,0.8)--cycle;
\node[align=center] at (0,0.2) {AStock推断\\概率、验证阈值、年末适用性};
\end{tikzpicture}
\end{exhibitbox}

事实可以直接陈述；预测必须注明来源与年份；推断必须说明假设和失效条件。下游需求、客户关系、平台认证、名义产能与股权绑定都不能独立证明订单、收入、利用率或利润。

\section{仍然没有解决的缺口}

截至截止日，未取得六只标的统一口径的当日官方北向持股、两融与席位数据，因此不以聚合数据替代。没有取得中国船舶、江波龙、盛新锂能与天华新能的显式当前外部目标；没有取得雅化与恩捷目标价的明确期限；恩捷客户涨价、实际单平利润与有效利用率仍未披露。

这些缺口不是报告“没做完”，而是公开证据边界。正确处理方式是零权、降级与持续验证，而不是填入未经证实的数字。

\clearpage

\appendix
\chapter{附录A：高空间假阳性审计}

\section{控股平台与投资收益}

吉林敖东、辽宁成大等公司在部分模型中可能出现极高目标空间，来源往往是把投资收益或持股价值折价修复，与高成长经营PE混合。投资收益受被投企业利润、市场价格与会计处理影响，不具有可重复的单位经济。正确方法是资产净值折价、分红能力与持股价值敏感性，而不是25--45倍成长PE。

\section{一次性或低基数利润}

九安医疗等标的在特殊年份会产生巨大现金与利润，若用单期利润年化，目标价可以轻易翻倍。研究必须区分持续经营利润、一次性项目、资产处置、政府补助与投资收益。没有可重复业务桥，就只能做资产与现金价值分析。

\section{周期峰值}

江波龙是本轮最典型的峰值分母风险。Q1毛利率55.53\%，H1利润预告92--110亿元，但现金流显著为负，券商2027利润增长又很低。若存储价格回落或补库成本上升，当前看似便宜的2026PE会迅速失真。

锂盐标的也有类似问题。锂价、低价库存、自供矿与套保可能在某一季度共振；只有持续出货、正常化吨利、现金与资产负债表一起改善，才能给估值信用。

\section{拥挤成长}

生益科技与深南电路作为AI硬件对照，基本面可能保持强劲，但高价与高预期会显著抬高翻倍门槛。一个好行业不自动产生一倍空间。若市场已经完整资本化未来数年的增长，任何轻微低于预期都可能引发倍数收缩。

\begin{exhibitbox}[Exhibit A1：假阳性类型与修复方法]
\small
\begin{tabularx}{\exhibitboxwidth}{L{3cm} L{3cm} X X}
\toprule
假阳性 & 典型标的 & 错误方法 & 修复方法 \\
\midrule
投资收益当成长EPS & 吉林敖东、辽宁成大 & 高成长PE & NAV折价、分红与资产价值 \\
一次性利润年化 & 九安医疗 & 单期利润乘倍数 & 正常化经营利润与现金价值 \\
周期峰值永续 & 江波龙、锂盐公司 & 峰值EPS乘高PE & 单位经济、现金流、周期中枢 \\
热门主题忽略价格 & 生益科技、深南电路 & 景气即翻倍 & 当前隐含增长与估值拥挤 \\
关系/认证当订单 & 多数成长公司 & 客户名录即收入 & 公司订单、转量、回款与利润 \\
\bottomrule
\end{tabularx}
\end{exhibitbox}

\clearpage

\chapter{附录B：模型台账与敏感性}

\section{翻倍所需倍数}

\begin{exhibitbox}[Exhibit B1：现价隐含与翻倍门槛]
\small
\begin{tabularx}{\exhibitboxwidth}{L{2.4cm} R{1.6cm} R{1.6cm} R{1.7cm} R{1.7cm} X}
\toprule
标的 & EPS分母 & 当前隐含PE & 翻倍所需PE & 牛市目标 & 解释 \\
\midrule
雅化集团 & 2026E 2.63 & 6.38x & 12.77x & 42.08 & 盈利兑现后仍需倍数约翻倍 \\
恩捷股份 & 2027E 5.16 & 9.27x & 18.54x & 103.20 & 接近完整20x恢复估值 \\
中国船舶 & 2027E 3.13 & 10.55x & 21.10x & 68.9 & 质量高但重估幅度大 \\
江波龙 & 2026E 26.47 & 14.68x & 29.35x & 780 & 峰值EPS上给高倍数 \\
盛新锂能 & 2026E 2.45 & 11.29x & 22.57x & 59.0 & 商品周期与现金双重验证 \\
天华新能 & 2026E 5.70 & 9.80x & 19.61x & 102.4 & 牛市目标低于翻倍价 \\
\bottomrule
\end{tabularx}
\end{exhibitbox}

\section{正式目标权重与同源敏感性}

只有雅化集团进入正式权重：80\%经营驱动概率值、10\%市场现价、10\%外部目标，悲观/基准/牛市概率为30\%/50\%/20\%。外部42元与基准经营参数来自同一份东吴报告，不能被视作第二个独立证据；10\%权重只用于保留可审计的Street参照。若将外部权重降为零并按比例重分至经营驱动与现价，结果为23.47元；若采用70\%经营驱动、20\%现价、10\%外部目标，结果为24.57元；主模型为25.32元。恩捷103元及30.96/67.08/103.20元三档全部为同源零权敏感性。

\section{情景不是预测分布}

概率并非通过历史样本统计得到，而是将证据强弱、目标期限、现金流与经营桥压缩成透明判断。它们的作用是避免牛市目标“偷渡”为基准结论。若新证据出现，概率、EPS和倍数必须全部更新。

\section{模型边界}

报告没有为零权备选发布正式目标，也没有把券商目标反向修改成年末目标。天华新能即使盈利强，牛市区间仍不翻倍；江波龙即使数学上刚好翻倍，证据与现金流不支持正式目标；中国船舶即使质量最好，缺失点目标与高倍数要求仍然阻止升级；恩捷虽有点目标，但公司级驱动证据阻断估值资格。

\clearpage

\chapter{附录C：主要来源与声明}

\section{公司公告}

主要官方材料包括六家公司2025年年度报告、2026年第一季度报告和2026年半年度业绩预告，以及恩捷股份限制性股票注销、中国船舶分红除息和江波龙限制性股票归属相关公告。公司预告均为未经审计估计，不能替代半年报。

\begin{itemize}
  \item 中国船舶（600150）：2025年报、2026Q1、2026H1业绩预告及分红实施股本公告；
  \item 江波龙（301308）：2025年报、2026Q1、2026H1业绩预告及股本变动公告；
  \item 恩捷股份（002812）：2025年报、2026Q1、2026H1业绩预告及限制性股票注销公告；
  \item 盛新锂能（002240）：2025年报、2026Q1、2026H1业绩预告；
  \item 天华新能（300390）：更正后的2025年报、2026Q1、2026H1业绩预告；
  \item 雅化集团（002497）：2025年报、2026Q1、2026H1业绩预告。
\end{itemize}

\section{原始券商报告}

本报告归档并阅读14份原始PDF。正式模型的限权外部锚只引用东吴证券《雅化集团：26H1业绩预告点评》（2026年7月7日，目标42元）；东吴证券《恩捷股份：26H1业绩预告点评》（2026年7月10日，目标103元）因公司级经营证据阻断，仅作零权敏感性。其他报告用于盈利预测、经营假设与风险比较，不用于点目标。

中国船舶覆盖来自五份公开原始报告；江波龙两份；盛新锂能四份；天华新能一份。公开集合全为正面评级，存在明显选择偏差。

\section{行情与市场环境}

个股价格使用2026年7月22日收盘后结构化行情，并以腾讯行情交叉核对。市场环境引用当日公开收盘报道：上证3,867.03点、深证成指下跌1.42\%、创业板下跌3.23\%、两市成交约2.65万亿元、超过3,800只股票下跌。

\section{使用声明}

本报告是截至2026年7月22日的研究快照，不构成证券买卖、仓位管理或收益承诺。目标价与概率均依赖公开信息和明确假设。市场、公司经营、会计口径与券商预测均可能变化；任何决策前应重新获取最新公告、行情与风险承受能力信息。

\begin{disclosurebox}[最终结论复述]
没有高置信度基准翻倍标的。雅化集团是唯一条件核心候选；恩捷股份虽有显式牛市目标，但因经营证据不足只保留零权敏感性；中国船舶是质量备选；其余标的保持周期或高风险观察。牛市目标不等于预期收益。
\end{disclosurebox}

\end{document}
